\scanpage{260}
\mspage{1}

\unsure{Catégories|Catégorie} de Chemins (nouvelle version)

Comme le formalisme ``simplifié'' des chemins que j'ai voulu dégager à la va-vite s'est révélé d'une lourdeur insupportable, de nature à le rendre quasiment inutilisable, je vais essayer de dégager un formalisme souple, collant à l'intuition, au prix seulement d'un effort conceptuel supplémentaire.

1. \uline{Intervalles} \uline{et \unsure{leur|leurs} \unsure{orientations|orientation}}.

On appelle \uline{catégorie} \uline{intervalle} \uline{\unsure{standard|standart}} une catégorie définie par \deleted{\unsure{un}} \unsure{ens.} ordonné, dont l'ens. des \deleted{\unsure{s}}éléments est un intervalle d'entiers $[0,n]$ ($n\in\mathbb{N}$, on n'exclut pas le cas $n=0$, i.e. de la catégorie ponctuelle $\{0\}$) :

\[\operatorname{Ob}I=[0,n]\subset\mathbb{N},\]

\deleted{\unsure{b)}} et satisfaisant l'hypothèse suivante : a) pour toute \deleted{\unsure{$i\in N$}, \unsure{l'un}} \inserted{paire de} \deleted{des sommets de $I$} deux sommets consécutifs $i-1,i\deleted{\unsure{+1}}$ de $I$ ($1\leq i\leq n\deleted{\unsure{+1}}$), il y a une flèche qui les relie, donc soit $i\longrightarrow i+1$, soit $i\longleftarrow i\unsure{+1|\text{et}}\deleted{\unsure{,\ \ldots}}$) Cette flèche est notée $\underline{u}_i$

\diagram{\begin{tikzcd}[ampersand replacement=\&,column sep=large]
(i-1)\arrow[r,"{\underline{u}_i}"] \& i \& \text{ou} \& (i-1) \& i\arrow[l,"{\underline{u}_i}"'] \& {;}
\end{tikzcd}}

b) La catégorie $I$ est engendrée par l'ens. de ses flèches $\underline{u}_i$ (appelées \uline{flèches de \unsure{transition}}.)

\deleted{Donc il y a au} La flèche \deleted{\unsure{élé}} \unsure{de transition} est dite \uline{directe} si elle va de $i-1$ à $i$, \uline{rétrograde} dans le sens inverse. Si on affecte l'intervalle $[i-1,i]$ du signe $+$ ou $-$ suivant que $\underline{u}_i$ est directe ou rétrograde, on voit donc que les

% archive sheet number: 259
